% ============================================================
%  The real thing: secp256k1 and Ethereum keys
% ============================================================

\section{The real setting: secp256k1 and Ethereum}

Ethereum (like Bitcoin) uses the curve \texttt{secp256k1} defined by the SECG standard:
\[
  E:\; y^2 = x^3 + 7 \quad \text{over } \mathbb{F}_p, \qquad p = 2^{256} - 2^{32} - 977 .
\]

\begin{center}
  \footnotesize
  \begin{tabular}{l l}
    \toprule
    Parameter & Value \\
    \midrule
    $a, b$ & $0,\ 7$ \\
    $p$ & \texttt{FFFFFFFF FFFFFFFF FFFFFFFF FFFFFFFF FFFFFFFF FFFFFFFF FFFFFFFE FFFFFC2F} \\
    $G_x$ & \texttt{79BE667E F9DCBBAC 55A06295 CE870B07 029BFCDB 2DCE28D9 59F2815B 16F81798} \\
    $G_y$ & \texttt{483ADA77 26A3C465 5DA4FBFC 0E1108A8 FD17B448 A6855419 9C47D08F FB10D4B8} \\
    $n$ (order of $G$) & \texttt{FFFFFFFF FFFFFFFF FFFFFFFF FFFFFFFE BAAEDCE6 AF48A03B BFD25E8C D0364141} \\
    $h$ (cofactor) & 1 \\
    \bottomrule
  \end{tabular}
\end{center}

The order $n$ is a 256-bit \textbf{prime} and the cofactor $h = N/n = 1$, so $G$ generates the whole curve. Compare with the toy curve, where $N = 18$, $n = 9$ (not prime) and $h = 2$.

\subsection{From private key to address}

\begin{figure}[H]
  \centering
  \begin{tikzpicture}[node distance=1.0cm, box/.style={draw, rounded corners=3pt, fill=gray!8, align=center, font=\small, minimum height=1.1cm, text width=3.0cm}]
    \node[box] (k) {Private key $k$\\ \scriptsize 256-bit random integer\\ \scriptsize $1 \le k < n$};
    \node[box, right=of k] (Q) {Public key $Q = kG$\\ \scriptsize point $(x, y)$,\\ \scriptsize 64 bytes $x \| y$};
    \node[box, right=of Q] (A) {Address\\ \scriptsize last 20 bytes of\\ \scriptsize $\mathrm{Keccak256}(x \| y)$};
    \draw[msg] (k) -- node[lbl, above] {ECC} node[lbl, below] {one-way} (Q);
    \draw[msg] (Q) -- node[lbl, above] {hash} node[lbl, below] {one-way} (A);
  \end{tikzpicture}
  \caption{Key derivation in Ethereum: two one-way steps of a different nature.}
  \label{fig:eth}
\end{figure}

\begin{enumerate}
  \item \textbf{Private key}: 32 random bytes, interpreted as an integer $k$, valid if $1 \le k \le n-1$. Its security depends entirely on the quality of the randomness.
  \item \textbf{Public key}: $Q = kG$ computed with scalar multiplication. Ethereum uses the uncompressed form: the 32-byte $x$ followed by the 32-byte $y$ (64 bytes, without the \texttt{0x04} prefix used by the SEC format).
  \item \textbf{Address}: Keccak-256 hash of the 64 bytes, keeping the last 20 bytes, written in hex with \texttt{0x} in front.
\end{enumerate}

\begin{framedobs}{Two different one-way functions in a row}
  The pipeline is a perfect illustration of the homework question. The first arrow (ECC) is one-way \textit{with structure}: the public key can be used to verify signatures made with $k$ (ECDSA). The second arrow (hash) is one-way \textit{without structure}: it just compresses the public key into a short, random-looking identifier. In fact, from a signature and the message, the public key can be recovered (\texttt{ecrecover}) and compared with the address.
\end{framedobs}
